% Corrected French for printed p.135.
% Corrections: an neau -> anneau; monomorphime -> monomorphisme;
% à axiome -> à l'axiome; note éditoriale sur la dernière assertion de 1.9.2.
\SourcePage{135}

\LayerTarget{tohoku:prop:1.9.2}
\textsc{Proposition 1.9.2.} \textit{Soit $\Sigma$ un schéma de
diagrammes avec relations de commutation \textup{(}\textit{cf} \textup{1.6)}, et soient $\mathbf{C}$ une
catégorie abélienne, $(U_i)_{i\in I}$ une famille de générateurs de $\mathbf{C}$. On
suppose que pour toute flèche de $\Sigma$, l'origine et l'extrémité de la
flèche sont distinctes, et que dans les relations de commutation ne figurent
pas les morphismes identiques $e_s$ (où $s$ est un sommet)\textsuperscript{2)}.
Alors pour tout $A\in \mathbf{C}$, et tout sommet $s$ du schéma, le diagramme
$\varepsilon_s(A)$ dont la valeur est $A$ au sommet $s$ et $0$ en tous les
autres sommets, et dont la valeur sur chaque flèche est $0$, appartient à
$\mathbf{C}^\Sigma$. De plus, la famille des $\varepsilon_s(U_i)$ ($s$ et $i$
variables) est un système de générateurs de $\mathbf{C}^\Sigma$.}

\LayerTarget{tohoku:prop:1.9.2:proof}
La vérification est immédiate, il suffit de remarquer, pour la dernière
assertion, que les morphismes de $\varepsilon_s(A)$ dans un diagramme $D$
s'identifient aux morphismes de $A$ dans $D(s)$.

\LayerTarget{tohoku:def:cogenerator-family}
Nous laissons au lecteur le soin de développer la notion duale de
\textit{famille de cogénérateurs} d'une catégorie abélienne. On peut montrer
que si $\mathbf{C}$ est une catégorie abélienne satisfaisant à l'axiome AB 5)
(\textit{cf.} 1.5.) alors l'existence d'un générateur implique l'existence d'un
cogénérateur (nous ne nous servirons pas de ce résultat). Ainsi la catégorie
des modules unitaires à gauche sur un anneau $U$ avec unité admet toujours un
cogénérateur : si par exemple $U=\mathbf Z$, on peut prendre pour cogénérateur
le groupe des nombres rationnels mod $\mathbf Z$ (ou le tore
$\mathbf T=\mathbf R/\mathbf Z$).

\LayerTarget{tohoku:section:1.10}
\textbf{1.10. Objets injectifs et projectifs.}
\LayerTarget{tohoku:def:injective-object}
Rappelons qu'un objet $M$ d'une catégorie abélienne $\mathbf{C}$ est dit
\textit{injectif} si le foncteur $A\to\mathrm{Hom}(A,M)$ (qui de toutes
façons est exact à gauche) est exact, i. e. si pour tout morphisme $u$ dans
$M$ d'un sous-truc $B$ d'un $A\in \mathbf{C}$, il existe un morphisme de $A$ dans $M$
qui le prolonge.
\LayerTarget{tohoku:def:injective-effacement}
Un morphisme $A\to M$ est dit un \textit{effacement injectif} (de
$A$) si c'est un monomorphisme, et si pour tout monomorphisme
$B\to C$ et tout morphisme $B\to A$, on peut trouver
un morphisme $C\to M$ rendant commutatif le diagramme

\LayerTarget{tohoku:diagram:injective-effacement}
\[
\begin{array}{ccc}
B & \longrightarrow & C \\
\big\downarrow && \big\downarrow \\
A & \longrightarrow & M
\end{array}
\]
Ainsi, pour que l'application identique de $M$ soit un effacement injectif il
faut et il suffit que $M$ soit injectif ; tout monomorphisme dans un objet
injectif est un effacement injectif.

\LayerTarget{tohoku:thm:1.10.1:p1}
\textsc{Théorème 1.10.1.} \textit{Si $\mathbf{C}$ satisfait à l'axiome
\textup{AB 5) (}\textit{cf.} \textup{1.5.)} et admet un générateur \textup{(}\textit{cf.} \textup{1.9)} alors pour tout $A\in \mathbf{C}$, il
existe un monomorphisme de $A$ dans un objet injectif $M$.}

On va même construire un foncteur
$M:A\to M(A)$ (non additif en général !)

\LayerTarget{tohoku:footnote:2}
\SourceFootnote{2}{La catégorie $\mathbf{C}^\Sigma$ des systèmes inductifs dans $\mathbf{C}$ construits sur un
ensemble d'indices \textit{ordonné} $I$, rentre dans ce cas. En effet il suffit
alors, dans l'exemple 1.7.h) de considérer les flèches $(i,j)$ avec $i<j$, et
les relations de commutation $(i,j)(j,k)=(i,k)$, où les $e_s$ n'interviennent
plus.}

\TohokuContinuation{136}
